% ==============================================================================
% CHƯƠNG 3: CHUẨN BỊ DỮ LIỆU (DATA PREPARATION)
% ==============================================================================

\chapter{Chuẩn Bị Dữ Liệu và Kiến Trúc Medallion (Data Preparation)}
\label{chap:data_preparation}

\section{Chiến Lược Thu Thập Dữ Liệu 3 Tầng (Three-Tier Pipeline)}

Để vừa bắt kịp các công trình nghiên cứu mới công bố hàng ngày, vừa đảm bảo chiều sâu lịch sử cho mạng lưới trích dẫn, hệ thống xây dựng quy trình thu thập dữ liệu tự động gồm 3 tầng (Three-Tier Harvesting Architecture) được lập lịch và điều phối bởi **Apache Airflow 2.10** theo múi giờ chuẩn Việt Nam (\texttt{Asia/Ho\_Chi\_Minh - ICT}):

\begin{enumerate}[label=\textbf{\arabic*.}]
  \item \textbf{Tier 1: Real-Time Stream (Hàng ngày lúc 02:00 AM ICT)}:
    \begin{itemize}
      \item \textit{Nguồn \& Phương thức}: Quét tự động các RSS Feeds chính thức của arXiv trên 5 chuyên ngành: \texttt{cs.AI}, \texttt{cs.LG}, \texttt{cs.CV}, \texttt{cs.CL}, và \texttt{stat.ML}.
      \item \textit{Trách nhiệm}: Bắt kịp các ấn phẩm tiền công bố trong vòng 24 giờ qua. Các bản ghi thu thập được đẩy ngay lập tức vào hàng đợi thông điệp RabbitMQ (\texttt{scientific\_paper\_ingest}) để worker bóc tách và nhúng vector vào Qdrant, phục vụ các truy vấn thời sự của người dùng.
    \end{itemize}
  \item \textbf{Tier 2: Near-Past Harvest (Hàng ngày lúc 02:10 AM ICT)}:
    \begin{itemize}
      \item \textit{Nguồn \& Phương thức}: Khai thác giao thức arXiv OAI-PMH với cơ chế phân trang \texttt{resumptionToken} và giới hạn tốc độ (\textit{rate-limiting delay} 6 giây/yêu cầu) nhằm tuân thủ quy định truy cập của arXiv.
      \item \textit{Trách nhiệm}: Thu thập khối lượng lớn **10,000 bài báo** từ các năm gần nhất (2024--2026), tự động chuyển đổi thành các phân vùng Parquet nén \texttt{zstd} trên MinIO và tạo điểm phục hồi dữ liệu tự động.
    \end{itemize}
  \item \textbf{Tier 3: Historical Landmark Archive (Định kỳ theo quý / Thủ công)}:
    \begin{itemize}
      \item \textit{Nguồn \& Phương thức}: Kết nối API OpenAlex, Semantic Scholar và Zenodo.
      \item \textit{Trách nhiệm}: Thu nạp các công trình nền tảng kinh điển (Foundation Papers như \textit{Attention Is All You Need}, \textit{Deep Residual Learning for Image Recognition - ResNet}, \textit{BERT}, \textit{Adam: A Method for Stochastic Optimization}). Đây là các đỉnh có bậc vào (in-degree) và điểm PageRank cực cao, đóng vai trò "xương sống" cho mạng lưới tri thức.
    \end{itemize}
\end{enumerate}

\section{Bóc Tách Văn Bản Học Thuật và Bảo Toàn Công Thức Toán}

Khác với các tài liệu văn phòng thông thường, tài liệu khoa học chứa đựng cấu trúc ngữ nghĩa chặt chẽ và hàm lượng phương trình toán học dày đặc. Module \texttt{AcademicHTMLParser} (triển khai tại \texttt{src/transformation/html\_parser.py}) được thiết kế chuyên biệt để phân tích các tài liệu HTML5 học thuật:

\begin{enumerate}[label=\alph*)]
  \item \textbf{Nhận dạng và phân loại Section}: Áp dụng các biểu thức chính quy (Regular Expressions) để phân loại từng đề mục vào các loại chuẩn: \texttt{abstract}, \texttt{introduction}, \texttt{related\_work}, \texttt{methodology}, \texttt{experiments}, \texttt{discussion}, \texttt{conclusion}, và \texttt{appendix}.
  \item \textbf{Bảo toàn tuyệt đối ký hiệu toán LaTeX}: Công cụ duyệt tìm các phần tử toán học dạng thẻ \texttt{<math>} hoặc \texttt{<span class="ltx\_Math">}, ưu tiên trích xuất chuỗi LaTeX gốc từ thẻ con \texttt{annotation encoding="application/x-tex"} hoặc thuộc tính \texttt{alttext}. Mỗi công thức được đóng gói chính xác trong cú pháp \lstinline{$...$} (inline math) hoặc \lstinline{$$...$$} (display math).
  \item \textbf{Loại bỏ nhiễu và làm sạch cấu trúc}: Tự động loại bỏ các thẻ script, navigation bar, style, header/footer của trang web, chuẩn hóa các khoảng trắng thừa nhưng bảo tồn nguyên vẹn cấu trúc đoạn văn (\texttt{<p>}).
\end{enumerate}

\section{Chuẩn Hóa Dữ Liệu và Cơ Chế Toàn Vẹn SHA-256}

Để đảm bảo tính nhất quán và loại bỏ triệt để dữ liệu rác trước khi lưu trữ vào tầng Silver, quy trình chuẩn hóa thực hiện các bước:
\begin{itemize}
  \item \textbf{Chuẩn hóa định danh tài liệu}: Chuyển đổi mã bài báo về dạng chuẩn không chứa tiền tố URL (ví dụ: chuyển \texttt{http://arxiv.org/abs/2401.13216v2} thành \texttt{2401.13216}).
  \item \textbf{Khử trùng lặp phiên bản (Deduplication)}: Nhóm các bản ghi theo \texttt{paper\_id}. Nếu một bài báo có nhiều phiên bản (v1, v2, v3), hệ thống sử dụng thuật toán sắp xếp ưu tiên phiên bản có tổng số công thức toán học (\texttt{total\_math\_count}) và tổng số từ vựng (\texttt{total\_words}) lớn nhất.
  \item \textbf{Mã băm toàn vẹn nội dung (SHA-256 Content Fingerprint)}: Mỗi tài liệu được tạo một mã băm SHA-256 dựa trên tổ hợp \texttt{paper\_id + title + clean\_full\_text}. Mã băm này đóng vai trò như một chứng thực bất biến (provenance ledger), ngăn chặn việc ghi đè hoặc tính toán lại vector cho các bài báo không thay đổi nội dung.
\end{itemize}

\section{Kiến Trúc Hồ Dữ Liệu Medallion (Medallion Lakehouse)}

Hệ thống triển khai kiến trúc **Medallion Lakehouse** chuẩn công nghiệp, phân tách triệt để giữa tầng tính toán (Compute) và tầng lưu trữ (Storage) nhằm tối ưu hóa chi phí và khả năng mở rộng:

\begin{enumerate}[label=\textbf{\arabic*.}]
  \item \textbf{Bronze Zone (\path{s3://uth-scientific-lakehouse/bronze/papers/YYYY/})}:
    Lưu trữ nguyên bản các tệp JSON siêu dữ liệu thô và tệp HTML5 học thuật được cào từ arXiv/ar5iv. Dữ liệu tại tầng này mang tính bất biến (immutable), cho phép hệ thống tái xử lý từ đầu bất cứ khi nào cập nhật thuật toán bóc tách mới.
  \item \textbf{Silver Zone (\path{s3://uth-scientific-lakehouse/silver/papers/year=YYYY/})}:
    Dữ liệu sau khi làm sạch được chuyển đổi sang định dạng cột **Apache Parquet**, nén bằng thuật toán \texttt{zstd} (Zstandard) mức nén tối ưu. Dữ liệu được phân vùng (partitioned) theo năm xuất bản. Tại tầng này, công cụ **DuckDB OLAP** có thể truy vấn trực tiếp dữ liệu từ S3 với kỹ thuật \textit{Predicate Pushdown} và \textit{Column Projection}, đạt tốc độ quét hàng trăm nghìn dòng chỉ trong vài mili-giây mà không cần nạp toàn bộ vào RAM.
  \item \textbf{Gold Zone}:
    \begin{itemize}
      \item \textit{Vector Store (Qdrant)}: Lưu trữ các vector biểu diễn ngữ nghĩa 768 chiều của các đoạn văn bản học thuật (\texttt{scientific\_papers\_gold}) được chỉ mục bằng đồ thị HNSW (Hierarchical Navigable Small World) viết bằng Rust, hỗ trợ lọc metadata thời gian thực theo chuyên ngành.
      \item \textit{Mining Artifacts (\path{s3://.../gold/mining/})}: Lưu trữ các tệp JSON tri thức cô đọng được xuất ra từ 4 Trụ cột Khai phá: đồ thị mạng lưới \texttt{graph\_coauthorship.json}, tập luật kết hợp \texttt{association\_rules.json}, phân cụm \texttt{clusters.json}, và danh sách dị biệt \texttt{trends\_anomalies.json}.
    \end{itemize}
\end{enumerate}

\section{Phân Đoạn Ngữ Cảnh Chuyên Sâu (Domain-Aware Deep Chunking)}

Hiệu năng của hệ thống tìm kiếm vector RAG phụ thuộc trực tiếp vào chiến lược phân đoạn văn bản. Nếu cắt văn bản theo số ký tự cố định một cách cơ học (fixed-character chunking), một đoạn trích có thể bị cắt ngang giữa một công thức toán học hoặc mất hoàn toàn thông tin ngữ cảnh rằng đoạn đó thuộc về phần nào của bài báo.

Module \texttt{AcademicChunker} (triển khai tại \texttt{src/indexing/chunker.py}) thực thi chiến lược phân đoạn sâu theo ngữ cảnh:

\begin{enumerate}[label=\alph*)]
  \item \textbf{Abstract Chunk độc lập}: Tóm tắt của bài báo luôn được đóng gói thành một chunk riêng biệt có mã \texttt{\{paper\_id\}\_abs}. Đây là chunk có giá trị ngữ nghĩa đại diện cao nhất cho toàn bộ công trình.
  \item \textbf{Phân đoạn theo ranh giới mục (Section-Aware Sliding Window)}: Duyệt qua từng section có cấu trúc trong \texttt{sections\_json}. Nếu một mục có độ dài vượt quá ngưỡng kích thước quy định ($\text{max\_chunk\_words} = 300\text{ từ}$), công cụ áp dụng cửa sổ trượt với độ chồng lấp $\text{overlap\_words} = 50\text{ từ}$.
  \item \textbf{Bổ sung ngữ cảnh nguồn (Contextual Enrichment)}: Mỗi đoạn trích sau khi cắt được tự động bổ sung tiêu đề bài báo và tên mục vào phần đầu:
  \begin{lstlisting}[language=Python]
context_text = f'Paper: "{title}" | Section: {sec_title}\n\n{chunk_text}'
  \end{lstlisting}
  Kỹ thuật này đảm bảo rằng khi một đoạn văn bản ngắn được nhúng thành vector, nó vẫn mang đầy đủ ngữ cảnh của bài báo chủ quản, ngăn chặn hiện tượng mất phương hướng ngữ nghĩa khi truy xuất độc lập.
\end{enumerate}

\section{Mô Hình Nhúng Vector Nomic Embed Text v1.5}

Hệ thống sử dụng mô hình nhúng mã nguồn mở **\texttt{nomic-ai/nomic-embed-text-v1.5}** với các đặc tính kỹ thuật vượt trội:
\begin{itemize}
  \item \textbf{Số chiều biểu diễn}: 768 chiều (768-D), cung cấp không gian hình học ngữ nghĩa phong phú gấp đôi so với các mô hình 384-D thế hệ cũ (như \texttt{all-MiniLM-L6-v2}).
  \item \textbf{Hỗ trợ ngữ cảnh dài}: Khả năng xử lý lên tới **8,192 tokens** nhờ cấu trúc mã hóa vị trí quay (Rotary Position Embeddings - RoPE), cấu hình trong hệ thống giới hạn an toàn ở mức 2,048 tokens bao phủ trọn vẹn mọi đoạn trích khoa học dài kèm công thức.
  \item \textbf{Cơ chế tiền tố bất đối xứng (Asymmetric Task Prefixing)}: Mô hình được huấn luyện đặc thù cho bài toán truy xuất thông tin không đối xứng:
    \begin{itemize}
      \item Khi nhúng câu hỏi của người dùng: Thêm tiền tố \texttt{search\_query: <câu hỏi>}.
      \item Khi nhúng đoạn tài liệu học thuật: Thêm tiền tố \texttt{search\_document: <văn bản>}.
    \end{itemize}
  \item \textbf{Chuẩn hóa L2 (L2 Normalization)}: Toàn bộ vector đầu ra được chuẩn hóa $\|v\|_2 = 1$, biến phép tính khoảng cách Cosine thành tích vô hướng đơn giản: $\text{CosineSim}(u, v) = u \cdot v$, tối ưu hóa tốc độ tìm kiếm trên đồ thị HNSW của Qdrant.
\end{itemize}
